\documentclass[12pt]{article}
\usepackage{amsmath, amssymb}

\begin{document}

\title{Lecture Scribe: Inequalities and Tail Bounds}
\author{Course: CSE400 – Fundamentals of Probability in Computing \\ Instructor: Dhaval Patel}
\date{}
\maketitle

\section*{1. Introduction and Motivation}

\subsection*{1.1 Definition of Tail}

The tail of a random variable $X$ is defined as:

\[
P(X > x)
\]

This represents the probability that the random variable takes values larger than a threshold $x$.

\subsection*{1.2 Why Do We Care About Tails?}

From the lecture:

\begin{itemize}
\item Jobs delayed more than 24 hours
\item Packets dropped due to full buffer
\end{itemize}

These represent rare but important events.

\subsection*{1.3 Key Idea}

Mean $\rightarrow$ average behavior (easy to compute) \\
Tail $\rightarrow$ rare extreme events (hard to compute)

\subsection*{1.4 Why Are Tails Hard to Compute?}

\begin{itemize}
\item Requires summing many small probabilities
\item Leads to long and complicated expressions
\end{itemize}

\section*{2. Examples of Tail Events}

\subsection*{2.1 Example 1: Load Balancing}

\textbf{Problem:} \\
Suppose $n$ jobs are distributed among $n$ servers at random. \\
Find:

\[
P(\text{a particular server gets } \geq k \text{ jobs})
\]

\subsubsection*{2.1.1 Intuition}

\begin{itemize}
\item Jobs are assigned randomly
\item Most servers get few jobs
\item Occasionally, one server gets many jobs
\end{itemize}

This is a tail event (rare overload).

\subsubsection*{2.1.2 Mathematical Formulation}

Let:

\[
X \sim \text{Binomial}(n, p)
\]

Then:

\[
P(X \geq k) = \sum_{i = k}^{n} \binom{n}{i} p^i (1 - p)^{n - i}
\]

\subsection*{2.2 Example 2: Poisson Arrivals}

\textbf{Problem:} \\
Jobs arrive according to a Poisson process with rate $\lambda$. \\
Find:

\[
P(\text{at least } k \text{ arrivals in one hour})
\]

\subsubsection*{2.2.1 Intuition}

\begin{itemize}
\item Jobs arrive randomly over time
\item Most intervals have moderate arrivals
\item Occasionally, bursts occur
\end{itemize}

This is a tail event (sudden spike).

\subsubsection*{2.2.2 Mathematical Formulation}

Let:

\[
X \sim \text{Poisson}(\lambda)
\]

Then:

\[
P(X \geq k) = \sum_{i = k}^{\infty} e^{-\lambda} \frac{\lambda^i}{i!}
\]

\section*{3. Why Tail Bounds Are Needed}

\subsection*{3.1 Exact Computation is Difficult}

From the lecture:

\[
P(X \geq k) = \sum_{i = k}^{\infty} e^{-\lambda} \frac{\lambda^i}{i!}
\]

\[
P(X \geq k) = \sum_{i = k}^{n} \binom{n}{i} p^i (1 - p)^{n - i}
\]

These involve:

\begin{itemize}
\item Long summations
\item Complex calculations
\end{itemize}

Hence:

Exact tail probabilities are often hard to compute

\subsection*{3.2 Core Idea of Tail Bounds}

Instead of computing exact values:

\begin{itemize}
\item Compute an upper bound
\item Easier to evaluate
\item Still provides a guarantee
\end{itemize}

Formally:

\[
P(X \geq k) \leq \text{simple expression}
\]

\subsection*{3.3 Key Insight}

Approximate safely instead of computing exactly

\section*{4. Running Example (Used Throughout Lecture)}

\subsection*{4.1 Setup}

Flip a fair coin $n$ times \\
Let:

\[
X = \text{number of heads}
\]

Then:

\[
X \sim \text{Binomial}\left(n, \frac{1}{2}\right)
\]

\subsection*{4.2 Mean}

\[
E[X] = \frac{n}{2}
\]

\subsection*{4.3 Tail Question}

\[
P\left(X \geq \frac{3n}{4}\right)
\]

This represents:

\begin{itemize}
\item Getting unusually many heads
\item A tail event
\end{itemize}

\subsection*{4.4 Purpose of Example}

\begin{itemize}
\item Develop progressively tighter bounds
\item Apply all inequalities on the same example
\end{itemize}

\section*{5. Markov’s Inequality}

\subsection*{5.1 Key Idea}

Uses only the mean \\
Large values are rare

Example from lecture:

If average $= 10$, very few values can be $\geq 100$

\subsection*{5.2 Statement of Markov’s Inequality}

For a non-negative random variable $X$:

\[
P(X \geq a) \leq \frac{E[X]}{a}
\]

\subsection*{5.3 Interpretation}

\begin{itemize}
\item Provides a simple upper bound
\item Requires only expectation
\item No need for full distribution
\end{itemize}

\section*{Conclusion (Covered Portion)}

From the lecture content covered:

\begin{itemize}
\item Tail probabilities measure rare extreme events
\item Exact computation is difficult due to long summations
\item Tail bounds provide simple, computable guarantees
\item Markov’s inequality is the first and simplest bound, using only the mean
\end{itemize}

\end{document}